\scanpage{22}
\uline{Démonstration} Comme $C\deleted{\unsure{B}}B\mathbin{\text{---}}CK$ et $GB\mathbin{\text{---}}GK$ sont surjectifs, on est ramené au cas de $B$.
Soit $\unsure{\xi^i|\xi}\in G^i(B)$ provenant de $P\in B^{(i)}$, on a
\[
\psi_B(\xi^i)=C_B^i(P\deleted{\unsure{(\{\lambda^i(X_s)\})}})
\]
d'où
\[
\varphi_B\psi_B(\xi^i)=\varphi_B(C_B^i(P))=c_B^i(P)
\equiv\deleted{\unsure{c_B^i}\ \text{\unsure{classe de}}}\gamma^i(P-\varepsilon(P))\bmod B^{(i+1)}.
\]
Or \unsure{on veut prouver}
\[
\gamma^i(P-\varepsilon(P))-(-1)^{i-1}(i-1)!P\in B^{(i+1)}
\]
\unsure{revient encore}, \uline{par} \uline{définition} de la filtration $B^{\deleted{\unsure{(i+1)}}}$, que l'image par $\widetilde C_B$ du premier mbr est dans $\widetilde{CB}^{(i+1)}$. Or cette image est
\[
\gamma^i(\overline P-\varepsilon(\overline P))-(-1)^{i-1}(i-1)!\overline P
\]
où $\overline P=\widetilde C_B(P)\in\deleted{\unsure{B^{(i)}}}\widetilde{CB}^{(i)}$.
Il est bien connu que \unsure{pour un tel $\overline P$}, l'expression précédente est dans $\widetilde{CB}^{(i+1)}$. On a bien prouvé
\[
\varphi_B\psi_B=(-1)^{i-1}(i-1)!\quad\text{en degré }i.
\]
On en conclut que
\[
(\psi_B\varphi_B)\psi_B=(-1)^{i-1}(i-1)!\psi_B\quad\text{en degré }i.
\]
Or l'\deleted{\unsure{espace}}\inserted{groupe-}but $CB$ de $\psi_B\varphi_B$ est sans \uline{torsion}, et le groupe image \deleted{de} $\psi_B$ dans $CB$ est $CB$ mod torsion, comme on vérifie \deleted{\unsure{très}} facilement.
\inserted{[en se ramenant au cas où $S$ est réduit à un élément ---]}
Il en résulte bien que $\psi_B\varphi_B=(-1)^{i-1}(i-1)!$, cqfd.

On voit que $\varphi_B:CB\longrightarrow GB$ est surjectif (\unsure{le cas se ramène au cas} où $S$ est réduit à un élément) donc \deleted{bijectif puisque} \inserted{$\varphi_B\psi_B=(-1)^{i-1}(i-1)!$}, et que $CB$ est sans torsion. Donc

\uline{Prop. 3} $\varphi_B$ est bijectif.

\uline{Corollaire} $\varphi_K$ est tjrs surjectif. Son noyau est un groupe de Torsion.

La première assertion résulte de ce que $GB\mathbin{\text{---}}GK$ est surjectif. La deuxième assertion résulte de $\psi_K\varphi_K=(-1)^{i-1}(i-1)!$.

\uline{Question} Est-ce que $\varphi_K$ est tjrs \uline{bijectif} ??? \uuline{Non}

\uline{Exemple} Soit \deleted{\unsure{par exemple}} $K=\mathbb Z[L]/(L-1)^n\deleted{\unsure{K}}$ ($\lambda_t L=1+Lt$), \inserted{(i.e. $K=K(P^{n-1})$)}, donc
\margin{\inserted{\uuline{un seul} générateur $X_s=L$}}
\margin{si $n-1>1$ i.e. $n\geq3$,}
\[
CK=\mathbb Z[\xi]/\text{\unsure{idéal homogène engendré par}}
\]
\[
(1+\xi)^{*n}=1+(-1)^{\unsure{n-1}}(n-1)!\xi^n+\text{---}
\]
et $CK$ a de la torsion. Cependant, \deleted{\unsure{on vérifie}} la filtration de $K$ est définie par les monômes isobares en
\[
\gamma^i(L-\varepsilon(L))=\gamma^i(L-1)
\]
qui sont \deleted{\unsure{tous}} égaux à: $L-1$ si $i=1$; à $0$ si $i>1$. Donc la filtration de $K$ est celle par les puissances de l'idéal \inserted{\uline{d'augmentation}} $(L-1)$, donc le gradué associé est sans torsion.

\margin{Mais si on adjoint à la liste géométrique $nL-1$ (qui correspond au fibré tangent à $\unsure{P(n-1)|P^{n-1}}$) on trouve $CK=\mathbb Z[\xi]/\xi^n$!}

\deleted{\uline{Proposition} (à prouver directement): Soit $K$ un $\lambda$-anneau, $(X_s)_{s\in S}$ une famille d'éléments géométriques de $K$, $K'$ le sous-$\lambda$-anneau engendré par les $X_s$, $C$ un \unsure{$\lambda$-homomorphisme} $K\longrightarrow\widetilde A$. Conditions équivalentes (i) $C^i(X_s)=0$ si $i>\varepsilon(X_s)$ (ii)}
\margin{21}
